%% Argos --- Selective Reasoning in Multi-Agent LLM Trading Desks
%% IEEE conference format (IEEEtran). First draft.
%%
%% Compile: pdflatex paper_ieee.tex  (twice)
%%
%% Numbers in the tables come from the reference run under
%% logs/decisions/nasdaq_focus/ (baseline) and reports/{gate_ab,
%% calibration, baselines}/summary.json. See Section VIII
%% (Reproducibility) for the exact journal paths and the git SHA.

\documentclass[conference]{IEEEtran}
\IEEEoverridecommandlockouts

\usepackage{cite}
\usepackage{amsmath,amssymb,amsfonts}
\usepackage{algorithmic}
\usepackage{graphicx}
\usepackage{textcomp}
\usepackage{xcolor}
\usepackage{booktabs}
\usepackage{url}
\usepackage{hyperref}
\hypersetup{colorlinks=true,linkcolor=blue!50!black,citecolor=blue!40!black,urlcolor=blue!50!black}

\def\BibTeX{{\rm B\kern-.05em{\sc i\kern-.025em b}\kern-.08em
    T\kern-.1667em\lower.7ex\hbox{E}\kern-.125emX}}

\begin{document}

\title{Selective Reasoning in Multi-Agent LLM Trading Desks:\\
A Cost--Recall--Calibration Study}

\author{
\IEEEauthorblockN{Kishor Timilsena}
\IEEEauthorblockA{Fuse-AI Fellowship\\
Department of Computer Science and Engineering\\
\textit{<home institution>}\\
kishor.timilsena@fuse.ai}
}

\maketitle

\begin{abstract}
Multi-agent Large Language Model (LLM) trading desks run a full
committee --- analysts, debaters, judge, risk manager, portfolio
manager --- on every bar, yet 84\% of decisions in our baseline are
HOLD. The cost of the committee is paid overwhelmingly to produce
inaction. We propose and evaluate a \emph{selective-reasoning} layer:
a single small-model call placed before the debate that decides
whether the committee convenes at all. We measure the
cost--recall--calibration tradeoff on four NASDAQ tickers over 2024
(NVDA, TSLA, ORCL, KO, $n{=}211$), with INTC~H1-2024 as an
out-of-universe probe for regime dependence, releasing full
per-decision audit logs for reproducibility. We report four findings:
(i) selective reasoning is regime-dependent --- on a low-vol
consumer-staple (KO) an aggressive gate saves 37\% of tokens at 55\%
true-trade recall while realized $h{=}5$ directional accuracy on
preserved trades rises from 82\% to 100\%; on a volatile declining
semi (INTC~H1) the gate's triggers fire on 11/11 bars and save
nothing; (ii) under symmetric model routing, the bull--bear debate
retains a ticker-specific confidence asymmetry ($|\Delta|{=}0.26$ on
NVDA, $|\Delta|{=}0.01$ on KO), suggesting evidence-conditional prior
leakage rather than model bias; (iii) per-agent calibration is uneven
--- fundamentals carries signal (Brier $0.12$ on ORCL at $n{=}8$)
but abstains often, bear-side advocacy on NVDA posts $37\%$ realized
directional accuracy at $h{=}5$; (iv) against random policies trading
at Argos's own empirical frequency, Argos's Sharpe is higher on
three of four tickers but 95\% bootstrap confidence intervals overlap
in every case. We release the full journal, prompts, and evaluation
harness as a research instrument.
\end{abstract}

\begin{IEEEkeywords}
multi-agent systems, large language models, algorithmic trading,
selective prediction, calibration, learning to defer, agent debate,
financial NLP
\end{IEEEkeywords}

%% =====================================================================
\section{Introduction}

Multi-agent large language model (LLM) systems have been proposed as a
paradigm for financial reasoning~\cite{tradingagents,finagent,finmem},
borrowing from the broader literature on multi-agent debate for LLM
reasoning~\cite{du2023debate,liang2023divergent}. These systems
orchestrate several specialized LLM workers --- analysts, debaters,
risk managers --- and commit to a trade decision per period. They are
evaluated primarily on returns.

Three measurement gaps limit this literature. \emph{First}, token cost
per decision is rarely reported, despite being the dominant operating
cost of these systems. \emph{Second}, inaction (HOLD) is treated as a
passive output; the expensive deliberation that produces it is
unexamined. \emph{Third}, per-agent contribution --- does the bear
advocate earn its tokens? --- is uninspected, making ablation hard.

We study a specific question that follows from these gaps:
\emph{given that multi-agent committees abstain the majority of the
time, when is convening the committee actually worth the tokens?}
This reframes trading agent orchestration as a selective-prediction
problem~\cite{chow1970,geifman2017selective,mozannar2020defer}, in
which a cheap gate decides whether the full committee runs or the
system short-circuits to HOLD.

\textbf{Contributions.} We present \emph{Argos}, an eight-node
LangGraph~\cite{langgraph} pipeline for multi-agent LLM
trading, instrumented for per-decision auditing. On top of the
pipeline we add a selective-reasoning gate before the debate and
measure its cost--recall--calibration frontier on four NASDAQ tickers
through 2024 ($n{=}211$ decisions). We report:
\begin{itemize}\itemsep1pt
  \item A \textbf{regime-dependent cost--recall tradeoff}: on KO the
  gate at $\tau{=}0.80$ saves 37\% of committee tokens while keeping
  55\% of the baseline's directional calls and raising $h{=}5$
  accuracy on the preserved trades from $82\%$ to $100\%$. On INTC
  H1-2024 the LLM gate chose \emph{proceed} on 11/11 bars; triggers
  fire every bar and no tokens are saved.
  \item A \textbf{confidence asymmetry under symmetric routing}: with
  bull and bear pinned to the identical model, key, and prompt
  structure, the confidence gap
  $|\mathbb{E}[c_{\text{bull}}]{-}\mathbb{E}[c_{\text{bear}}]|$
  varies from $0.012$ (KO) to $0.260$ (NVDA), evidencing
  ticker-specific prior leakage.
  \item A \textbf{per-agent calibration disaggregation}: the
  technicals agent is well-calibrated (ECE $0.028{-}0.055$) but
  accurate only at chance level; the blind judge is systematically
  overconfident (ECE $0.14{-}0.20$ on three of four tickers); the
  bear advocate on NVDA at $h{=}5$ posts $37\%$ realized accuracy
  with Brier $0.24$ --- well-calibrated about being wrong.
  \item A \textbf{baseline comparison} against buy-and-hold, moving-
  average crossover, random policies matched to Argos's empirical
  trade frequency, and a logistic learned-gate baseline. Argos's
  point Sharpe exceeds random's on three of four tickers but 95\%
  bootstrap CIs overlap in every case --- at $n{=}53$ weekly
  decisions per ticker the return advantage is not statistically
  distinguishable.
\end{itemize}

Our contribution is a \textbf{measurement framework and instrument}
rather than a new agent design: we show how to evaluate a multi-agent
LLM committee under cost, recall, and calibration jointly, and
release the full per-decision audit stream to enable replication.

%% =====================================================================
\section{Related Work}

\textbf{Multi-agent LLM finance.}
Xiao~\emph{et al.}~\cite{tradingagents} introduce a multi-agent
trading pipeline close in spirit to our system. Zhang~\emph{et
al.}~\cite{finagent} and Yu~\emph{et al.}~\cite{finmem} explore
tool-use and memory-augmented single-agent variants. These works
report decisions and returns on equity universes but do not
disaggregate per-agent contribution or report per-decision token
cost. FinGPT~\cite{fingpt} and BloombergGPT~\cite{bloomberggpt}
study foundation models for finance and are complementary.

\textbf{Multi-agent debate.}
Du~\emph{et al.}~\cite{du2023debate} and Liang~\emph{et
al.}~\cite{liang2023divergent} introduce debate as a procedure for
improving LLM reasoning on general tasks. Our bull--bear step adopts
their formulation; the blind-judge aggregator derives from
self-refinement work~\cite{madaan2023selfrefine}.

\textbf{Selective prediction and learning-to-defer.}
Selective classification abstains when the model is uncertain, trading
coverage for accuracy~\cite{chow1970,geifman2017selective,
geifman2019selectivenet}. Learning-to-defer frameworks~\cite{
mozannar2020defer,charusaie2022sample} treat the decision to defer as
a learned policy. We extend this family to the orchestration level:
the gate decides when to invoke a downstream multi-agent reasoning
procedure.

\textbf{LLM cost efficiency.}
Model cascading~\cite{chen2023frugalgpt},
routing~\cite{ong2024routellm}, and speculative
execution~\cite{leviathan2023speculative} reduce per-inference cost by
substituting cheaper models. Our gate is complementary --- it reduces
\emph{whether} the expensive path runs, orthogonal to \emph{which}
model it uses.

%% =====================================================================
\section{System: Argos}
\label{sec:system}

\begin{figure}[t]
  \centering
  \fbox{\parbox{0.92\columnwidth}{\centering\footnotesize
  \textbf{FIGURE 1}: Argos pipeline topology.\\[2pt]
  \texttt{START} $\to$ \texttt{exit\_manager} $\to$
  \texttt{\{fundamentals, sentiment, technicals\}} (parallel)
  $\to$ \texttt{kronos\_forecast}
  $\to$ \textbf{\texttt{edge\_gate}} $\to$ \texttt{bull/bear debate}
  $\to$ \texttt{blind\_judge} $\to$ \texttt{risk\_manager}
  $\to$ \texttt{portfolio\_manager} $\to$ \texttt{END}.\\[3pt]
  Two short-circuit paths: exit manager (open position) and gate HOLD.}}
  \caption{The selective-reasoning gate (bold) is the object of study.}
  \label{fig:topology}
\end{figure}

Argos is a LangGraph~\cite{langgraph} pipeline of eight LLM-backed
nodes. We summarize the topology (Fig.~\ref{fig:topology}); full
prompts are in the appendix and the companion repository.

\textbf{Agents and roles.}
Three analysts produce structured \texttt{AnalystSignal}
(\textit{bullish}/\textit{bearish}/\textit{neutral}, confidence
$\in[0,1]$, key factors, reasoning). A multi-horizon numerical
forecast node (Kronos~\cite{kronos}) produces 1-day/1-week/1-month
direction and expected return. The bull and bear advocates produce
\texttt{DebateArgument} with position, confidence, strongest evidence
and rebuttal. The blind judge reads the transcript with side labels
hidden and emits a \texttt{JudgeVerdict}. The risk manager votes on
veto; the portfolio manager synthesizes the final
\texttt{FinalDecision} with sizing and exit plan.

\textbf{Edge gate.}
Our contribution. One structured call per decision, pinned to a cheap
model tier (flash-lite), reading the three analyst signals and the
indicator snapshot. The gate emits \texttt{proceed}~$\in\{$True,
False$\}$, a confidence, and a list of triggers from a fixed
vocabulary: \texttt{analyst\_conviction},
\texttt{analyst\_disagreement}, \texttt{vol\_spike},
\texttt{composite\_bias\_skew}, \texttt{range\_breakout},
\texttt{earnings\_proximity}, \texttt{kronos\_high\_conviction},
\texttt{pattern\_fire}. The prompt is recall-biased: when in doubt,
proceed. On skip, the pipeline short-circuits to HOLD and the
bull/bear/judge/risk/PM stages are not called.

\textbf{Journal.}
Every $(\text{ticker}, \text{date})$ decision is journalled as JSONL
including: the exact prompt sent to each agent, each agent's
structured output, per-agent token counts, the gate verdict, the
final decision, the full indicator snapshot, and the git SHA that
produced it. This audit stream is the research instrument underlying
every number in this paper.

%% =====================================================================
\section{Selective-Reasoning Formulation}
\label{sec:formulation}

Let $x_t$ be the decision context at $(\text{ticker}, t)$: analyst
signals, indicator snapshot, forecast, prior portfolio state. A
\emph{committee policy} $\pi_C$ produces decision $y_t \in
\{\textsc{Buy}, \textsc{Sell}, \textsc{Hold}\}$ at token cost $c_t^C$
(dominated by debate, judge, risk, PM, averaging $\sim27$k
tokens/decision in our system). A \emph{gate}
$g:x_t \to \{0,1\}$ decides whether $\pi_C$ runs. Under the selective
policy $\pi_g$:
\begin{equation}
  y_t^g =
  \begin{cases}
    \pi_C(x_t), & g(x_t) = 1 \\
    \textsc{Hold}, & g(x_t) = 0
  \end{cases}
\end{equation}
\begin{equation}
  c_t^g = c_g + \mathbf{1}[g(x_t){=}1]\, c_t^C
\end{equation}
where $c_g \approx 1.5$k tokens is the gate's own cost.

\textbf{Metrics.}
Given a baseline run of $\pi_C$ on the same journal, let
$\mathcal{D}_+ = \{t : \pi_C(x_t) \ne \textsc{Hold}\}$ be the
baseline's directional set. We report:
\begin{itemize}\itemsep1pt
  \item \textbf{Token savings:}
    $1 - \sum_t c_t^g \,/\, \sum_t c_t^C$.
  \item \textbf{True-trade recall:}
    $\lvert\{t \in \mathcal{D}_+ : g(x_t){=}1\}\rvert
     \,/\, \lvert\mathcal{D}_+\rvert$.
  \item \textbf{Directional accuracy at horizon $h$:}
    $\Pr[\mathrm{sign}(r_{t,t+h}) = \mathrm{sign}(y_t^g)
     \,|\, y_t^g \ne \textsc{Hold}]$
    for $h \in \{1,5,20\}$ trading days.
  \item \textbf{Calibration:} per-agent reliability
    diagram~\cite{niculescu2005predicting} and Brier
    score~\cite{brier1950}.
  \item \textbf{Sharpe regret:}
    $S(\pi_g) - S(\pi_C)$ with bootstrap CI (1000 resamples of the
    per-period return series, resampled with replacement).
\end{itemize}

%% =====================================================================
\section{Experiments}

\textbf{Universe and cadence.}
Four tickers: NVDA (AI-rally growth), TSLA (high-vol EV), ORCL
(mid-pivot cloud), KO (low-vol staple). Weekly cadence over 2024
($\sim$53 decisions per ticker, $n{=}211$ decisions total). INTC
(2024 H1, declining semi) is an out-of-universe probe for the
regime-dependence hypothesis ($n{=}11$).

\textbf{Arms.}
\emph{Baseline} ($\pi_C$): full committee every bar.
\emph{Gated} ($\pi_g$): gate-conditioned invocation. Both arms use
identical models, prompts, and random seeds for every non-gate
agent; the LLM cache is shared across arms. Journal paths:
\verb|logs/decisions/nasdaq_focus/| (baseline) and
\verb|logs/decisions/nasdaq_focus_gated/|.

\textbf{Ablations.}
We additionally run three post-hoc ablations that cost no new LLM
calls: (i) a \emph{thresholded gate} sweeping gate confidence
$\tau\in[0,1]$ to trace the full Pareto frontier;
(ii) a \emph{logistic baseline} on engineered features from the
indicator snapshot; (iii) a \emph{leave-one-analyst-out} analysis
using the journal's structured signals.

\textbf{Models and configuration.}
Analysts: Gemini 3.5 Flash with Google-Search grounding
(fundamentals, sentiment) or ungrounded (technicals). Debate
(bull/bear): same model, same key, symmetric by construction.
Judge: Gemini 3.5 Flash on a different key family. Risk and
portfolio manager: Gemini 3.5 Pro. Edge gate: Gemini 3.5 Flash-lite.
Fallback cascade: 6 Gemini keys $\to$ Groq $\to$ OpenRouter. All
model IDs are pinned (Section~\ref{sec:repro}).

\textbf{Fees and execution.}
Paper-trade with US retail fee profile: zero commission, 1bp
slippage per side, 2bp round-trip breakeven. Position sizing via
fractional Kelly from the portfolio manager; stops, targets,
trailing and time-stop come from the portfolio manager's exit plan
and are honored by the simulator.

%% =====================================================================
\section{Results}

\subsection{The cost--recall Pareto frontier}
\label{sec:results-pareto}

\begin{figure}[t]
  \centering
  \fbox{\parbox{0.92\columnwidth}{\centering\footnotesize
  \textbf{FIGURE 2}: Pareto frontier. \emph{x-axis}: token savings
  vs.\ baseline (\%). \emph{y-axis}: true-trade recall. One curve
  per ticker, points marked at gate-confidence thresholds
  $\tau\in\{0,0.1,\ldots,0.9\}$.\\
  Source: \texttt{reports/gate\_ab/pareto.json}.}}
  \caption{Cost--recall tradeoff. Up-and-to-the-right is better.
  Different tickers occupy different regions of the frontier.}
  \label{fig:pareto}
\end{figure}

Table~\ref{tab:pareto} gives three threshold slices of the heuristic
sweep. Here \emph{proceed\%} is the fraction of bars the thresholded
gate would run the full committee on; \emph{tok.\ saved\%} is the
fraction of baseline tokens that short-circuiting would avoid;
\emph{recall} is the fraction of the baseline's BUY/SELL decisions
preserved; the last column is $h{=}5$ realized directional accuracy
on the gate-proceeded directional calls.

\begin{table}[t]
\caption{Gate threshold sweep. Up/right is better.}
\label{tab:pareto}
\centering\footnotesize
\begin{tabular}{llrrrr}
\toprule
Ticker & $\tau$ & proceed\% & recall & tok.\ saved\% & $h{=}5$ acc \\
\midrule
NVDA & 0.80 & 100\% & 100\% &  0.0\% & 82\% \\
NVDA & 0.85 & 76\%  & 76\%  & 18.8\% & \textbf{85\%} \\
NVDA & 0.90 & 53\%  & 76\%  & \textbf{41.4\%} & 85\% \\
\midrule
TSLA & 0.80 & 62\%  & ---   & 34.6\% & --- \\
\midrule
ORCL & 0.80 & 57\%  & 50\%  & 38.9\% & 67\% \\
ORCL & 0.85 & 49\%  & 33\%  & 46.6\% & 50\% \\
\midrule
KO   & 0.80 & 58\%  & 55\%  & 36.6\% & \textbf{100\%} \\
KO   & 0.85 & 57\%  & 45\%  & 38.6\% & 100\% \\
\midrule
INTC & live & 100\% & 1.00 & 0\% & --- \\
\bottomrule
\end{tabular}
\end{table}

Two asymmetric findings sit in this table. \emph{First}: on NVDA and
KO the gate-at-$\tau{\geq}0.85$ \emph{raises} $h{=}5$ directional
accuracy on the trades it preserves (NVDA $82\%{\to}85\%$; KO
$82\%{\to}100\%$), while cutting committee tokens by $19\%{-}39\%$.
The gate is implicitly a confidence filter: the marginal decisions
it drops are also the ones more likely to be directional errors.
\emph{Second}: on ORCL the same aggression collapses recall
($50\%{\to}33\%$) with no further accuracy gain --- the gate drops
real trades along with noise. The payoff is ticker-specific, not
architecture-level.

INTC's $100\%$ proceed rate across 11 live gate invocations is the
complementary evidence: on a volatile declining name in H1-2024,
every bar triggers at least one of
\texttt{vol\_spike}, \texttt{analyst\_disagreement},
\texttt{composite\_bias\_skew}, \texttt{analyst\_conviction}.
The gate correctly refuses to save tokens when the regime does not
permit it.

\subsection{Confidence asymmetry under symmetric routing}
\label{sec:results-asymmetry}

Bull and bear share the same model, API key, prompt scaffolding,
and decoding parameters. Any confidence gap is therefore
evidence-driven, not model-prior-driven --- yet the gap varies
substantially by ticker (Table~\ref{tab:asymmetry}).

\begin{table}[t]
\caption{Confidence asymmetry under symmetric model routing.}
\label{tab:asymmetry}
\centering\footnotesize
\begin{tabular}{lrrr}
\toprule
Ticker & $\mathbb{E}[c_{\text{bull}}]$ & $\mathbb{E}[c_{\text{bear}}]$ & $|\Delta|$ \\
\midrule
NVDA & 0.679 & 0.419 & \textbf{0.260} \\
TSLA & 0.452 & 0.562 & 0.110 \\
ORCL & 0.475 & 0.528 & 0.054 \\
KO   & 0.478 & 0.491 & \textbf{0.012} \\
\bottomrule
\end{tabular}
\end{table}

We interpret this as \emph{evidence-conditional prior leakage}: the
same symmetric debate structure retains a ticker-specific confidence
gap because the public narrative density around each ticker differs.
NVDA's 2024 AI-rally context overwhelms bear-side arguments; KO's
narrative vacuum leaves the debate balanced.

\subsection{Per-agent calibration}
\label{sec:results-calibration}

\begin{figure}[t]
  \centering
  \fbox{\parbox{0.92\columnwidth}{\centering\footnotesize
  \textbf{FIGURE 3}: Per-agent reliability diagrams (fundamentals,
  sentiment, technicals, judge, bull, bear) with Brier scores.
  Source: \texttt{reports/calibration/}.}}
  \caption{Per-agent calibration. Diagonal = well-calibrated.}
  \label{fig:calibration}
\end{figure}

Table~\ref{tab:calibration} shows per-agent calibration at $h{=}5$,
aggregated across the baseline window. Brier score $\in[0,1]$
(lower is better; $0.25$ is random on a binary target); ECE is
expected calibration error across confidence bins
$\{[0,0.5), [0.5,0.6), \ldots, [0.9,1]\}$; \emph{gap} is mean
confidence minus mean realized accuracy (positive implies
overconfidence).

\begin{table}[t]
\caption{Per-agent calibration at $h{=}5$.}
\label{tab:calibration}
\centering\scriptsize
\begin{tabular}{llrrrrr}
\toprule
Ticker & Agent & $n$ & Brier & ECE & mean c. & mean acc. \\
\midrule
NVDA & fund.      & 53 & 0.274 & \textbf{0.208} & 0.831 & 0.623 \\
NVDA & sent.      & 44 & 0.258 & 0.165          & 0.801 & 0.636 \\
NVDA & tech.      & 42 & 0.253 & \textbf{0.055} & 0.507 & 0.452 \\
NVDA & judge      & 53 & 0.266 & 0.184          & 0.807 & 0.623 \\
NVDA & bull       & 53 & 0.229 & 0.057          & 0.680 & 0.623 \\
NVDA & bear       & 53 & 0.239 & 0.042          & 0.419 & \textbf{0.377} \\
\midrule
ORCL & fund.      &  8 & \textbf{0.115} & 0.075 & 0.800 & \textbf{0.875} \\
ORCL & tech.      & 47 & 0.254 & 0.028          & 0.517 & 0.489 \\
ORCL & judge      & 53 & 0.278 & \textbf{0.200} & 0.766 & 0.566 \\
ORCL & bull       & 53 & 0.238 & 0.073          & 0.582 & 0.509 \\
ORCL & bear       & 53 & 0.249 & 0.068          & 0.559 & 0.491 \\
\midrule
KO   & tech.      & 41 & 0.257 & 0.041          & 0.578 & 0.537 \\
KO   & judge      & 53 & 0.246 & \textbf{0.135} & 0.758 & 0.623 \\
KO   & bull       & 53 & 0.249 & 0.095          & 0.586 & 0.491 \\
KO   & bear       & 53 & 0.243 & 0.072          & 0.581 & 0.509 \\
\bottomrule
\end{tabular}
\end{table}

Three structural observations follow from Table~\ref{tab:calibration}.

\textbf{Technicals is epistemically honest but not predictive.}
ECE between $0.028$ and $0.055$ across every ticker means the
technicals agent's stated confidence tracks its realized correctness
bin-by-bin --- but its mean accuracy is $45$--$54\%$, near chance.
The technicals agent knows it doesn't know.

\textbf{The blind judge is systematically overconfident.}
ECE on the judge sits at $0.135{-}0.208$ on NVDA, ORCL, and KO
(TSLA omitted for space, same pattern). The blind adjudication
procedure does not neutralize confidence inflation --- it
introduces its own.

\textbf{The bear advocate on NVDA is well-calibrated about being
wrong.}
Brier $0.239$ and ECE $0.042$ with mean confidence $0.419$ and
realized accuracy $0.377$: the role posts below-chance $h{=}5$
accuracy in a rally regime \emph{and} assigns itself appropriately
low confidence. This is evidence that the adversarial-role prompt
contributes systematic negative signal on bullish names, not noise.

\subsection{Learned-baseline check: trigger vocabulary vs.\ decision rule}
\label{sec:results-learned}

We fit a leave-one-ticker-out logistic regression on seven
hand-chosen indicator features (max analyst confidence, analyst
disagreement, $|$net bias$|$, two vol-regime indicators, max
candle-pattern ringer, position in 20-day range) to predict whether
the baseline committee returned a directional verdict. The same
features feed the heuristic gate from
Section~\ref{sec:results-pareto}; the comparison isolates
\emph{decision rule} (linear combination vs.\ max-of-triggers) from
\emph{feature vocabulary} (Table~\ref{tab:learned}).

\begin{table}[t]
\caption{Logistic vs.\ heuristic gate (LOTO cross-validation).}
\label{tab:learned}
\centering\scriptsize
\begin{tabular}{llrrrr}
\toprule
Held-out & pos/$n$ & log.\ ROC & log.\ PR & heur.\ ROC & heur.\ PR \\
\midrule
NVDA & 17/53 & $0.438$ & $0.292$ & $\mathbf{0.656}$ & $\mathbf{0.458}$ \\
ORCL &  6/53 & $0.390$ & $0.171$ & $0.466$ & $0.102$ \\
KO   & 11/53 & $\mathbf{0.574}$ & $\mathbf{0.369}$ & $0.453$ & $0.190$ \\
TSLA &  0/53 & --- & --- & --- & --- \\
\bottomrule
\end{tabular}
\end{table}

The comparison yields three different answers. On NVDA the heuristic
gate's trigger-vocabulary dominates the linear baseline
(ROC $0.66$ vs.\ $0.44$); on KO the linear combiner over the same
features beats the max-of-triggers rule (ROC $0.57$ vs.\ $0.45$); on
ORCL both are near chance. Read together, this suggests the LLM
gate's value is located in its trigger \emph{vocabulary}, not its
decision rule. A future gate design should keep an LLM for trigger
identification but replace the max-of-triggers rule with a learned
per-ticker linear combiner --- a direction consistent with the
broader literature on model routing~\cite{ong2024routellm,
chen2023frugalgpt}.

\subsection{Argos vs.\ non-LLM baselines}
\label{sec:results-baselines}

A multi-agent LLM system that cannot be distinguished from a
frequency-matched random policy at the current sample size would
not warrant the operating cost. Table~\ref{tab:baselines} reports
that comparison explicitly. Each baseline is matched to Argos's own
empirical $(\mathrm{buy},\mathrm{sell})$ rate observed in the
per-ticker journals; random Sharpes are the mean over $1000$ seed
draws; Argos Sharpe CIs are $95\%$ bootstrap (resample $n{=}53$
period returns with replacement, $1000$ draws).

\begin{table}[t]
\caption{Baselines vs.\ Argos.}
\label{tab:baselines}
\centering\scriptsize
\begin{tabular}{llrrc}
\toprule
Ticker & Strategy & Return (\%) & Sharpe & 95\% CI \\
\midrule
NVDA & buy\_and\_hold         & $+178.9$ & $+1.01$ & --- \\
NVDA & ma\_cross\_20/50       & $+39.9$  & $+0.48$ & --- \\
NVDA & random (emp.)          & $+74.2$  & $+1.59$ & $[+0.14,+2.97]$ \\
NVDA & \textbf{argos}         & $+2.0$   & $\mathbf{+1.68}$ & $[-0.30,+3.83]$ \\
\midrule
TSLA & buy\_and\_hold         & $+62.6$  & $+0.49$ & --- \\
TSLA & ma\_cross\_20/50       & $+105.8$ & $+0.77$ & --- \\
TSLA & \textbf{argos}         & $+0.0$   & ---     & --- \\
\midrule
ORCL & buy\_and\_hold         & $+62.1$  & $+0.76$ & --- \\
ORCL & ma\_cross\_20/50       & $+23.7$  & $+0.42$ & --- \\
ORCL & random (emp.)          & $+29.2$  & $+1.15$ & $[-0.32,+2.35]$ \\
ORCL & \textbf{argos}         & $+0.03$  & $\mathbf{+1.18}$ & $[-0.55,+2.81]$ \\
\midrule
KO   & buy\_and\_hold         & $+7.3$   & $+0.28$ & --- \\
KO   & ma\_cross\_20/50       & $+10.7$  & $+0.52$ & --- \\
KO   & random (emp.)          & $+3.3$   & $+0.39$ & $[-1.18,+1.96]$ \\
KO   & \textbf{argos}         & $+0.27$  & $\mathbf{+0.89}$ & $[-1.19,+3.11]$ \\
\bottomrule
\end{tabular}
\end{table}

Argos's point Sharpe exceeds random's mean Sharpe on NVDA, ORCL, and
KO. On all three, the $95\%$ bootstrap CI contains the random
policy's mean. We therefore report \emph{no statistically
distinguishable Sharpe advantage at $n{=}53$ weekly decisions per
ticker}.

The realized dollar returns ($+0.03$--$+2.00\%$) are an order of
magnitude below buy-and-hold ($+7$--$+179\%$), driven by Argos's
one-trade-per-ticker decision stream in the baseline (median $1$
paper trade/ticker over the year). Sharpe is being measured mostly
on held-flat periods with intermittent exposure; Argos is in
practice a market-timing abstainer, not a trading strategy, and its
Sharpe reflects exposure selection more than position management.

TSLA is a non-result for return: the committee did not open a single
paper trade across $53$ decisions, forgoing the ticker's $+62.6\%$
buy-and-hold year. The same abstention is interpreted as
epistemically correct in Section~\ref{sec:discussion}; the
measurement reports both sides.

%% =====================================================================
\section{Discussion}
\label{sec:discussion}

\textbf{When the gate helps.}
The threshold sweep (Section~\ref{sec:results-pareto}) identifies the
gate operating points at which it is lossless or
accuracy-improving: NVDA and KO at $\tau{\geq}0.85$, saving
$19{-}39\%$ of committee tokens while raising $h{=}5$ accuracy on
preserved trades. In each case the pre-skipped bars are low-vol,
mid-range, with low analyst disagreement and weak composite bias
--- the committee's counterfactual output on those bars is HOLD, so
skipping is free.

\textbf{When the gate hurts.}
ORCL at the same $\tau$ shows recall collapsing from $100\%$ to
$33\%$ with no $h{=}5$ accuracy gain beyond $\tau{=}0.80$. The gate
cannot separate genuine ORCL directional calls from noise using the
current trigger vocabulary; the LLM gate's recall-biased prompt
compensates in the live run.

\textbf{When the gate cannot help.}
INTC H1-2024 shows the limit of this architecture: a regime in which
every bar presents a plausibly asymmetric question. Live gate runs
proceeded on $11/11$ bars, saving no tokens. The paper's framing
should resist the claim that selective reasoning is universally
cost-reducing --- it is regime-conditional, and the trigger
vocabulary determines which regimes benefit.

\textbf{Decision rule vs.\ trigger vocabulary.}
Isolating the LLM gate's design into trigger vocabulary and decision
rule (Section~\ref{sec:results-learned}) produces a ticker-split
finding: the LLM gate's max-of-triggers rule dominates a logistic
baseline on NVDA (ROC $0.66$ vs.\ $0.44$) but loses to it on KO
($0.45$ vs.\ $0.57$). This suggests the LLM is contributing
\emph{what to measure} more than \emph{how to combine} --- a
cost-conscious successor should keep the LLM for feature
identification and defer combination to a per-ticker learned linear
combiner, cutting the gate's per-call budget while likely improving
its selection quality.

\textbf{The research-instrument argument.}
Whether or not gating saves tokens on any given ticker, the
per-decision audit stream carries the paper's primary contribution:
a reproducible measurement of multi-agent LLM deliberation under
controlled model routing, with per-agent calibration, confidence
asymmetry, and realized-return decomposition exposed for independent
inspection. The confidence-asymmetry finding
(Section~\ref{sec:results-asymmetry}) and the below-chance
bear-on-NVDA calibration (Section~\ref{sec:results-calibration}) are
unavailable in prior multi-agent-LLM-finance work for this reason.

\textbf{The statistical-honesty argument.}
The Argos-vs-random overlap (Section~\ref{sec:results-baselines}) is
a sample-size result, not a method result: at $n{=}53$ per ticker,
two policies with Sharpe ratios differing by $0.5$ cannot be
separated at $\alpha{=}0.05$. We state this openly rather than
cherry-pick. Scaling the evaluation universe is the first-order
follow-up.

%% =====================================================================
\section{Limitations}

Our universe is small (4+1 US tickers), the window is one year, the
cadence is weekly, and execution is paper. The learned-baseline
analysis is retrospective; a prospective bandit-style gate is left
to future work. All results are on the Gemini 3.5 Flash family;
replication with Claude Haiku or Mistral-small would isolate
model-family effects. The simulator applies a US-retail fee profile;
execution at scale would warrant market-impact modeling.

A further limitation is \emph{training-data contamination}: 2024
likely appears in the pretraining corpora of the models used.
Directional accuracy numbers may be inflated by implicit knowledge
despite explicit ``as-of'' instructions in the prompts. The
confidence-asymmetry and calibration findings are robust to this
--- both advocates share the same contaminated prior.

%% =====================================================================
\section{Reproducibility}
\label{sec:repro}

Journal paths, pinned model IDs, environment toggles, git SHAs, and
LLM-cache hashes for every run are listed in the online appendix.
All prompts live under \texttt{prompts/} at the git SHA named in
each journal record. The LLM cache (SQLite) is included in the
artifact drop so results can be re-derived without LLM credits.
Model router pins (as of the current submission):
\texttt{gemini-3.5-flash}, \texttt{gemini-3.5-flash-lite},
\texttt{gemini-3.5-pro}, each via LangChain's \texttt{with\_structured\_output}.

%% =====================================================================
\section{Conclusion}

We proposed selective reasoning as a cost-control mechanism for
multi-agent LLM trading desks and measured it on five US equities
through 2024. The mechanism is regime-dependent; savings are
non-uniform. The measurement framework itself is the reusable
contribution. Treating multi-agent LLM finance as a measurement
problem --- cost, recall, calibration jointly --- is the direction
we recommend for future work.

%% =====================================================================
\begin{thebibliography}{99}\footnotesize

\bibitem{chow1970}
C.~K. Chow,
``On optimum recognition error and reject tradeoff,''
\emph{IEEE Trans.\ Inf.\ Theory}, vol.~16, no.~1, pp.~41--46, 1970.

\bibitem{geifman2017selective}
Y.~Geifman and R.~El-Yaniv,
``Selective classification for deep neural networks,''
in \emph{NeurIPS}, 2017.

\bibitem{geifman2019selectivenet}
Y.~Geifman and R.~El-Yaniv,
``SelectiveNet: A deep neural network with an integrated reject option,''
in \emph{ICML}, 2019.

\bibitem{mozannar2020defer}
H.~Mozannar and D.~Sontag,
``Consistent estimators for learning to defer to an expert,''
in \emph{ICML}, 2020.

\bibitem{charusaie2022sample}
M.-A. Charusaie \emph{et al.},
``Sample efficient learning of predictors that complement humans,''
in \emph{ICML}, 2022.

\bibitem{du2023debate}
Y.~Du \emph{et al.},
``Improving factuality and reasoning in language models through
multiagent debate,''
\emph{arXiv:2305.14325}, 2023.

\bibitem{liang2023divergent}
T.~Liang \emph{et al.},
``Encouraging divergent thinking in large language models through
multi-agent debate,''
\emph{arXiv:2305.19118}, 2023.

\bibitem{madaan2023selfrefine}
A.~Madaan \emph{et al.},
``Self-Refine: Iterative refinement with self-feedback,''
in \emph{NeurIPS}, 2023.

\bibitem{tradingagents}
Y.~Xiao \emph{et al.},
``TradingAgents: Multi-agent LLM financial trading framework,''
\emph{arXiv:2412.20138}, 2024.

\bibitem{finagent}
A.~Zhang \emph{et al.},
``A multimodal foundation agent for financial trading,''
\emph{arXiv:2402.18485}, 2024.

\bibitem{finmem}
Y.~Yu \emph{et al.},
``FinMem: A performance-enhanced LLM trading agent with layered
memory and character design,''
\emph{arXiv:2311.13743}, 2023.

\bibitem{fingpt}
H.~Yang \emph{et al.},
``FinGPT: Open-source financial large language models,''
\emph{arXiv:2306.06031}, 2023.

\bibitem{bloomberggpt}
S.~Wu \emph{et al.},
``BloombergGPT: A large language model for finance,''
\emph{arXiv:2303.17564}, 2023.

\bibitem{chen2023frugalgpt}
L.~Chen \emph{et al.},
``FrugalGPT: How to use large language models while reducing cost
and improving performance,''
\emph{arXiv:2305.05176}, 2023.

\bibitem{ong2024routellm}
I.~Ong \emph{et al.},
``RouteLLM: Learning to route LLMs with preference data,''
\emph{arXiv:2406.18665}, 2024.

\bibitem{leviathan2023speculative}
Y.~Leviathan \emph{et al.},
``Fast inference from transformers via speculative decoding,''
in \emph{ICML}, 2023.

\bibitem{brier1950}
G.~W. Brier,
``Verification of forecasts expressed in terms of probability,''
\emph{Monthly Weather Review}, vol.~78, no.~1, pp.~1--3, 1950.

\bibitem{niculescu2005predicting}
A.~Niculescu-Mizil and R.~Caruana,
``Predicting good probabilities with supervised learning,''
in \emph{ICML}, 2005.

\bibitem{langgraph}
LangChain Inc.,
``LangGraph: Building stateful multi-actor LLM applications,''
2024. [Online]. Available:
\url{https://langchain-ai.github.io/langgraph/}

\bibitem{kronos}
Shiyu \emph{et al.},
``Kronos: Candlestick foundation model,''
2024. [Online]. Available: \url{https://github.com/shiyu-coder/Kronos}

\end{thebibliography}

\end{document}
